% Standalone problem document, generated from the adjacent README.md.
% Compile with XeLaTeX. Mathematical content is maintained in Markdown.
\documentclass[11pt,a4paper]{article}
\usepackage[fontsize=10.5pt]{scrextend}
\usepackage[margin=22mm,headheight=15pt,headsep=9mm,footskip=11mm]{geometry}
\usepackage{fontspec}
\setmainfont{texgyrepagella}[Extension=.otf,UprightFont=*-regular,BoldFont=*-bold,ItalicFont=*-italic,BoldItalicFont=*-bolditalic]
\setsansfont{texgyreheros}[Extension=.otf,UprightFont=*-regular,BoldFont=*-bold,ItalicFont=*-italic,BoldItalicFont=*-bolditalic]
\setmonofont{texgyrecursor}[Extension=.otf,UprightFont=*-regular,BoldFont=*-bold,ItalicFont=*-italic,BoldItalicFont=*-bolditalic,Scale=MatchLowercase]
\usepackage{amsmath,amssymb}
\usepackage{unicode-math}
\setmathfont{texgyrepagella-math.otf}
\usepackage{xcolor,microtype,enumitem,needspace,fancyhdr}
\usepackage{bookmark}
\definecolor{ink}{HTML}{17344B}
\definecolor{accent}{HTML}{197D80}
\definecolor{muted}{HTML}{596773}
\hypersetup{colorlinks=true,linkcolor=accent,urlcolor=accent,pdftitle={IE-08 - A
cubic-time Schur algorithm using logarithmic
precision},pdfauthor={Open Problems in Numerical Linear Algebra}}
\urlstyle{same}
\setlength{\parindent}{0pt}
\setlength{\parskip}{5pt plus 1pt minus 1pt}
\linespread{1.04}
\setlength{\emergencystretch}{3em}
\widowpenalty=10000
\clubpenalty=10000
\displaywidowpenalty=10000
\allowdisplaybreaks[1]
\setlist{nosep,leftmargin=1.5em}
\providecommand{\tightlist}{\setlength{\itemsep}{0pt}\setlength{\parskip}{0pt}}
\providecommand{\pandocbounded}[1]{#1}
\pagestyle{fancy}
\fancyhf{}
\fancyhead[L]{\sffamily\footnotesize\color{muted}OPEN PROBLEMS IN NUMERICAL LINEAR ALGEBRA}
\fancyhead[R]{\sffamily\bfseries\footnotesize\color{accent}IE-08}
\fancyfoot[L]{\parbox[b]{0.9\textwidth}{\sffamily\fontsize{8}{10}\selectfont\color{muted}Editorial ratings; a bounded search cannot certify that a problem remains unsolved.\\Literature check: 2026-09-11}}
\fancyfoot[R]{\sffamily\footnotesize\color{muted}\thepage}
\renewcommand{\headrulewidth}{0.3pt}
\renewcommand{\headrule}{\hbox to\headwidth{\color{accent}\leaders\hrule height \headrulewidth\hfill}}
\makeatletter
\renewcommand{\section}{\Needspace{5\baselineskip}\@startsection{section}{1}{0pt}{12pt plus 2pt minus 2pt}{4pt}{\sffamily\bfseries\large\color{ink}}}
\renewcommand{\subsection}{\Needspace{4\baselineskip}\@startsection{subsection}{2}{0pt}{9pt plus 1pt minus 1pt}{3pt}{\sffamily\bfseries\normalsize\color{ink}}}
\makeatother
\setcounter{secnumdepth}{0}
\begin{document}
{\sffamily\small\color{muted}Eigenvalues and inverse problems\par}
\vspace{3pt}
{\raggedright\hyphenpenalty=10000\exhyphenpenalty=10000\sffamily\bfseries\fontsize{21}{24}\selectfont\color{ink}A
cubic-time Schur algorithm using logarithmic precision\par}
\vspace{7pt}
{\sffamily\small\textbf{Difficulty:} extreme\quad\textbf{Importance:} broadly
interesting\par}
{\sffamily\small\bfseries\color{accent}Status: Solved\par}
\vspace{4pt}
{\color{accent}\hrule height 0.8pt}
\vspace{6pt}
\textbf{Rating rationale:} Extreme reflects an end-to-end precision and
complexity theorem for general Schur decomposition; broad impact is
justified by this central eigenvalue-computation task.

\subsection{Resolution --- 2026-09-11}\label{resolution-2026-09-11}

\textbf{Affirmative resolution by Matthew J. Colbrook}, Department of
Applied Mathematics and Theoretical Physics, University of Cambridge.
\href{https://github.com/ajt60gaibb/OpenProblemsInNLA/blob/main/eigenvalues-and-inverse-problems/IE-08/solution.md}{Complete
manuscript} ·
\href{https://github.com/ajt60gaibb/OpenProblemsInNLA/blob/main/eigenvalues-and-inverse-problems/IE-08/solution.pdf}{PDF}
·
\href{https://github.com/ajt60gaibb/OpenProblemsInNLA/blob/main/eigenvalues-and-inverse-problems/IE-08/solution.tex}{LaTeX}.
\textbf{Theorem 1, Lemmas 2--12 and the final four proof sections.}

The end-to-end algorithm uses \(O(n^3\log^c(n/\delta))\) arithmetic
operations and \(O(\log(n/\delta))\) mantissa bits, with universal
constants, for every complex input with \(\|A\|_2\le1\). With
probability at least \(0.99\) it returns an exactly upper triangular
\(T\) and a \(Q\) satisfying both displayed residual bounds. The proof
includes finite random sampling, deterministic work caps, global
recursive conditioning and floating-point error control, without an
input separation or diagonalizability assumption. It is an asymptotic
existence result with conservative constants, not a production
implementation.

The complete argument received an independent Codex-agent \textbf{PASS}
on 11 September 2026. The
\href{https://github.com/ajt60gaibb/OpenProblemsInNLA/blob/main/references/colbrook-additional-2026-09-11/verification/reviews/IE-08-review.md}{review
report} records the exact scope and a hash of the original reviewed
manuscript. The mathematical sections remain unchanged in the authored
version. Original ChatGPT generation is disclosed; no external human
peer review or formal proof certificate is asserted.
\href{https://github.com/ajt60gaibb/OpenProblemsInNLA/blob/main/references/colbrook-additional-2026-09-11/README.md}{Submission
and verification record}.

The original statement, source references and prior audit notes are
retained; the former difficulty rating is historical.

\subsection{Problem statement}\label{problem-statement}

Does a randomized floating-point algorithm exist with the following
guarantee? Given \(A\in\mathbb C^{n\times n}\), \(\|A\|_2\leq1\), and
\(0<\delta<1/2\), it uses at most

\[
O(n^3\log^c(n/\delta))\quad\text{arithmetic operations and}\quad
O(\log(n/\delta))\quad\text{mantissa bits}
\]

and, with probability at least \(0.99\), returns an upper triangular
\(T\) and \(Q\) satisfying

\[
\|A-QTQ^*\|_2\leq\delta,
\qquad \|Q^*Q-I\|_2\leq\delta?
\]

The constants and exponent \(c\) must be universal, with no additional
eigenvalue-separation or diagonalizability assumption. Use the usual
relative-error floating-point model, with sufficient exponent range to
avoid overflow and underflow. This is a statement about the precision
needed for a rigorous algorithm, not whether ordinary QR implementations
usually work well.

\newpage
\subsection{References}\label{references}

Amsel et al.,
\href{https://arxiv.org/html/2602.05394v3\#S3.SS2}{\emph{Linear Systems
and Eigenvalue Problems}}, Problem 3.3. Schneider,
\href{https://escholarship.org/content/qt3bb8s95w/qt3bb8s95w_noSplash_22c8cf7d2873d111b5bb367d09dc40fc.pdf}{\emph{Pseudospectral
Divide-and-Conquer for the Generalized Eigenvalue Problem}},
dissertation, §1.6.1 and Chapter 6. Banks et al.,
\href{https://doi.org/10.1007/s10208-022-09577-5}{FOCM diagonalization
paper}, §6, precision-reduction question.

\subsection{Earlier status check ---
2026-09-08}\label{earlier-status-check-2026-09-08}

Searches for \textit{Schur logarithmic precision 2026} and
\textit{site:arxiv.org Schur precision 2026} found no complete
analysis meeting both bounds. The dissertation explicitly distinguishes
analyzed subroutines from an end-to-end precision theorem.

\subsection{Audit update --- 2026-09-10}\label{audit-update-2026-09-10}

Rechecked \href{https://arxiv.org/html/2602.05394v3}{workshop version
3}, Problem 3.3 and its update. The report still distinguishes
algorithmic ingredients from a complete logarithmic-precision analysis.
Targeted cubic-Schur and precision searches found no theorem meeting
every displayed requirement.
\end{document}
